In a small controlled comparison at hand-set default weights, an energy-deviation running cost made iLQR-MPC swing-up more successful than the quadratic cost on the pendulum and cheaper in control effort on both plants. Most of that advantage is tied to the default quadratic weight. The success gap is absent on the cart-pole (ceiling) and disappears on the pendulum with a larger quadratic weight. The cart-pole effort gap shrinks from \rhval{main/ilqr-energy-ours/cartpole/control_effort/mean:2} vs \rhval{main/ilqr-quad/cartpole/control_effort/mean:2} at default weights to \rhval{tuned_all/ilqr-energy-ours/cartpole/control_effort/mean:1} vs \rhval{tuned_all/ilqr-quad/cartpole/control_effort/mean:1} when each cost uses its best swept weight; on the pendulum the tuned gap is \rhval{tuned_all/ilqr-energy-ours/pendulum/control_effort/mean:1} vs \rhval{tuned_all/ilqr-quad/pendulum/control_effort/mean:1}. A lower effort for the energy cost at its default or tuned weight is thus the one effect that kept its direction against every quadratic weight we tried. Its size is modest once the quadratic cost is re-weighted, it rests on 5 seeds, three of them shared with weight selection, and on the cart-pole it does not hold for every energy weight. The energy cost was no less sensitive to its weight, was indistinguishable from the combined cost on the pendulum, and did not beat energy shaping with LQR catch there (energy shaping used less effort). Testing the costs under model error, with jointly tuned weights on separate seeds and tuned baselines, is the natural next step.
